% ===========================================================================
\subsection{\file{new-design-reference/assets/ASSETS.md}}
% ===========================================================================

\textbf{Markdown, 33 lines.} A checklist of the twelve files somebody would need
to produce before this design could be considered finished, and three notes about
placeholder content to replace.

\begin{provenance}{new-design-reference/assets/ASSETS.md}
Hand-written, added in the first commit \code{0d77d51} alongside the design it
describes, and never edited since. Written by the same person in the same sitting
as the HTML, which is evident from the fact that it references the exact class
names (\code{pt1} through \code{pt6}, \code{pa1}, \code{pa2}) used in that file.
\textbf{Confidence: high, from the commit history and the internal references.}
\end{provenance}

\subsubsection*{Why it exists}

It is a handover note. The design was built with placeholder visuals, and this
file records precisely what would replace each one and what HTML change each
replacement needs. As an artefact it is genuinely good practice: it converts
``this is not finished'' from a feeling into a list.

\begin{jargon}{Markdown}
A way of writing formatted text in a plain file, using punctuation instead of
menus: \code{\#} starts a heading, backticks mark code, and pipe characters draw a
table. It is designed to be readable as plain text and is the standard format for
notes kept alongside code.
\end{jargon}

\subsubsection*{Line by line}

\begin{lstlisting}[firstnumber=1,caption={\file{ASSETS.md}, lines 1--21}]
# Assets checklist: ai.digitalise.agency

Drop real assets here, then update the HTML as noted.

## Required before launch

| File | Size | Where used | HTML change needed |
|---|---|---|---|
| `videos/hero-1.mp4` | under 8MB | Hero background tiles | Swap .tile-blue .tile-thumb div for `<video autoplay muted loop playsinline>` |
| `logos/client-1.svg` | any | Proof section | Replace `pa1`/`pa2` avatar initials with `<img>` |
| `logos/client-2.svg` | any | Proof section | Same |
| `logos/client-3.svg` | any | Proof section | Same |
| `thumbs/port-1.jpg` | 800x450px | Portfolio card 1 | Add `style="background-image:url('/assets/thumbs/port-1.jpg')"` to `.pt1` |
| `thumbs/port-2.jpg` | 800x450px | Portfolio card 2 | Same pattern for `.pt2` through `.pt6` |
| `thumbs/port-3.jpg` | 800x450px | Portfolio card 3 | |
| `thumbs/port-4.jpg` | 800x450px | Portfolio card 4 | |
| `thumbs/port-5.jpg` | 800x450px | Portfolio card 5 | |
| `thumbs/port-6.jpg` | 800x450px | Portfolio card 6 | |
| `favicon.svg` | 32x32 | Browser tab | Add `<link rel="icon" href="/assets/favicon.svg">` to `<head>` |
| `og-image.png` | 1200x630px | Social share preview | Update `og:image` meta tag |
\end{lstlisting}

Line 1 is the title, and line 3 the instruction: put the real files in this folder,
then change the HTML as the table says.

Lines 7 and 8 open a table, with line 8 being the separator row that Markdown
requires between a table's headings and its body.

The twelve rows are the substance. Reading them tells you exactly what the design
is pretending to have:

\begin{itemize}
\item \textbf{One video}, under 8 megabytes, to play inside the hero tiles instead
  of the static icons. The suggested tag carries \code{autoplay muted loop
  playsinline}, and the \code{muted} is not optional: browsers refuse to autoplay
  video with sound, so an unmuted autoplaying video simply does not start.
  \code{playsinline} stops iPhones taking the video full screen.
\item \textbf{Three client logos}, to replace the coloured circles with initials in
  the proof section. Note the row says \code{pa1}/\code{pa2} and there are three
  logos and three avatars (\code{pa1} to \code{pa3}), so the instruction is one
  short.
\item \textbf{Six portfolio thumbnails} at 800 by 450, which is exactly the 16 by 9
  shape the stylesheet reserves with \code{aspect-ratio} on line 473.
\item \textbf{A favicon}, with the exact tag needed to install it.
\item \textbf{A social preview image} at 1200 by 630, the standard Open Graph size,
  to fix the missing \code{og:image} noted in section 7.1.
\end{itemize}

\begin{honest}{None of the twelve exist}
The \file{assets/} folder contains exactly one file: this checklist. Every asset it
calls ``required before launch'' is absent, which means the design was never
finished to the standard its own author set for it. The portfolio thumbnails are
still gradients, the client logos are still lettered circles, and the page has
neither a favicon nor a social preview image.
\end{honest}

\begin{lstlisting}[firstnumber=21,caption={\file{ASSETS.md}, lines 21--33}]

## Numbers to update (index.html metrics band)

Search for `data-target` in `index.html` and replace all four values with real numbers verified from Gantry / n8n data.

## Booking link

Search for `YOUR_CAL_LINK` in `index.html` and replace with the actual Cal.com or Calendly URL.
For a full inline calendar embed, replace the `<a>` tag with the Cal.com embed snippet from cal.com/docs/embeds.

## Client proof section

Three `<!-- UPDATE -->` comments in the proof section mark where real client names and companies go.
\end{lstlisting}

Line 24 is the most interesting line in the file, because of one phrase: the four
metric numbers are to be replaced with ``real numbers verified from Gantry / n8n
data''. Gantry is another project in this workspace, the agency's content
publishing system, and n8n is the automation tool behind several of these projects.
So the author's intention was that the numbers on this marketing page would be
traceable to the agency's own operational systems rather than estimated. That is a
higher standard than most marketing pages hold themselves to, and it makes the
invented hero tile numbers in \file{index.html} more surprising rather than less.

Line 28 tells the reader to search for \code{YOUR\_CAL\_LINK}.

\begin{honest}{Two of this file's three instructions are stale}
Searching \file{index.html} for \code{YOUR\_CAL\_LINK} finds nothing. The literal
placeholder was replaced on 2026-07-28 by commit \code{4dffc12}, which set the link
to \code{cal.com/davonex/30min}. So somebody following this instruction searches,
finds nothing, and has no way to tell whether the job is done or the instruction is
wrong. The answer is that the job is half done: the placeholder text is gone and the
link still points at another project's calendar.

The metrics instruction on line 24 is stale in the same way: commit \code{2c66a7a}
claims to have replaced all placeholder content with real agency data, and the four
values do look like considered numbers rather than round placeholders. Yet the
\code{UPDATE} comment sits in \file{index.html} at line 187 and this instruction sits
here, both saying the work is outstanding.

Only the third instruction, on line 33 about the client proof section, is verifiably
still true: the avatar initial on line 430 of \file{index.html} is a leftover.
\end{honest}

\subsubsection*{What it connects to}

It describes \file{new-design-reference/index.html} and
\file{new-design-reference/css/site.css} by name and by class, and it is the only
file in the repository that documents the project's intent. Nothing reads it: it is
for people.

The deploy workflow explicitly excludes it from publication (line 82 of the
workflow), which is the right instinct, and that exclusion no longer matches its
path. See section 7.6.

\subsubsection*{How it breaks}

As a document it cannot break. As a source of truth it already has: two of its
three instructions describe a state of the code that no longer exists, which is the
characteristic failure of documentation kept beside code and not updated with it.

